\section{Ameaças à Validade}
\label{sec:ameacas}

A estrutura desta seção segue a taxonomia usual de ameaças em engenharia de software
experimental \cite{wohlin2012}. Cada ameaça é apoiada em uma evidência do próprio estudo e
acompanhada da mitigação adotada, ou da razão de não ter sido possível mitigá-la.

\textbf{Validade de Construto.} A equivalência entre release e deploy, e entre
\emph{workflow run} vermelho e falha de entrega, é uma aproximação inerente ao uso de
\emph{proxies}. A amostra-ouro mitiga parcialmente a primeira: 295 das 300 releases
rotuladas representam entregas reais ao usuário, 5 ficaram incertas e nenhuma foi rotulada
como "não" --- a suposição se sustenta nesta amostra. Ela permanece, porém, para as 95
releases de bibliotecas e \emph{frameworks}, em que publicar no registro não implica que
alguém atualizou a dependência. A segunda equivalência é o objeto da Seção
\ref{sec:discussao} e não foi mitigada, e sim \emph{medida}: as duas variantes do CFR
discordam em 72,3\% dos repositórios. A heurística corretiva que sustenta o CFR (b) tem
F1 de 0,754, com 55 falsos positivos e 27 falsos negativos, e superdetecta corretivas ---
marca 65,4\% das releases contra 51\% na rotulagem manual. Além disso,
\texttt{commit.author.date} é vulnerável a \emph{rebase} e \emph{squash merge}, que reescrevem
a data e distorcem o lead time sem relação com o esforço real; mitigamos apenas
contabilizando datas ausentes em vez de estimá-las, e mantendo valores negativos visíveis em
vez de descartá-los. Por fim, há uma ameaça ao próprio instrumento de validação: o consenso
da amostra-ouro foi consolidado em um arquivo único do grupo, e as planilhas individuais dos
três avaliadores não estão versionadas, o que impede recalcular o $\kappa$ de Fleiss e,
portanto, aferir a confiabilidade da rotulagem que sustenta o F1 relatado.

\textbf{Validade Interna.} O principal confundidor identificado é a própria cadência: um
repositório que publica com frequência tem mais oportunidades de ter uma release corretiva
dentro de sete dias, o que infla o CFR (b) independentemente da qualidade. Esse mecanismo
foi quantificado (87,7\% contra 25,0\% dos intervalos dentro da janela) e é reportado junto
do coeficiente, não omitido. Releases publicadas por robôs de automação ---
\emph{release-please}, \emph{semantic-release}, Dependabot --- entram como qualquer outra e
inflam a frequência de projetos que automatizaram a publicação; essa ameaça não foi mitigada,
pois não distinguimos autor humano de robô. Execuções agendadas são excluídas pelo filtro de
evento \texttt{push}, mas execuções com matriz de combinações contam como um único
\emph{run}. O truncamento no fim da janela foi tratado por censura explícita, e não por
descarte, em ambos os pontos em que ocorre: 196 episódios de recuperação (0,97\%) e 246
releases do CFR (b) (1,3\%).

\textbf{Validade Externa.} A amostra é deliberadamente estreita: repositórios públicos com
pelo menos 1.001 estrelas, que usam GitHub Actions, observados em uma única janela anual. O
funil dimensiona o recorte: 511 dos 1.021 candidatos avaliados caíram por publicar menos de
cinco releases em doze meses, e outros 103 por não usarem Actions. Os achados descrevem,
portanto, projetos populares, ativos e automatizados --- não o GitHub como um todo, e muito
menos desenvolvimento corporativo, onde o DORA foi formulado e onde "deploy" tem acesso a
dados de produção que aqui não existem. Pipelines em GitLab CI ou Jenkins estão fora do
alcance do instrumento. A janela única impede ainda qualquer afirmação de tendência: 2025 é
um retrato, não uma série temporal.

\textbf{Validade de Conclusão.} A família principal da RQ 06 executa 16 testes, e a correção
de Holm foi aplicada sobre ela; o efeito é observável --- estrelas contra frequência tem
$p = 0{,}0070$ bruto e $p = 0{,}0556$ corrigido, deixando de ser significativo. Todos os
testes são não paramétricos, o que protege as deduções contra a assimetria documentada na
Figura \ref{fig:violino}, e todos os tamanhos de efeito são reportados junto dos $p$-valores.
A análise exploratória por tipo de projeto apoia-se em 60 repositórios distribuídos em três
tipos, e está declarada como exploratória por isso. A mediana de recuperação usa apenas
episódios recuperados e, por construção, subestima o tempo real, já que os episódios mais
longos são justamente os que não terminam na janela; publicamos o limite inferior que trata
cada censurado como recuperado no fim da observação (5,22 h contra 4,94 h) para dimensionar
esse viés. O $n$ difere entre métricas (264 na recuperação contra 300 nas demais), o que
impede comparações diretas de dispersão entre elas. Por fim, as faixas DORA produzem muitos
empates de fronteira: 86,3\% dos repositórios estão a um degrau de outra categoria, de modo
que comparações de categoria entre estudos precisam reportar essa proximidade, sob pena de
atribuir a erro de cálculo o que é efeito da discretização.
